% related_work.tex -- drop into the ICLR 2027 main file after the Introduction.
% Requires: \bibliography{references}
%   (references.bib in this directory: the 74 new entries plus everything from papers/shared/references.bib,
%    so one \\bibliography{references} resolves every citation in the paper.)
% Style matches the Introduction: no em dashes, no colons or semicolons in running prose.

\section{Related Work}
\label{sec-related}

\paragraph{Sycophancy is a trained bias, and it is steerable.}
Preference-tuned models agree with the user more often than the evidence
warrants, and the deference grows with model scale and with the confidence
of the user \citep{perez2023discovering,sharma23}. It survives into
consequential settings, where a model validates the asker's own account of
a conflict \citep{cheng25elephant} and where it revises a correct answer
under pushback \citep{fanous2025syceval,malmqvist24}. Two facts about the
bias matter here. It is a linear direction in activation space that can be
added or subtracted \citep{rimsky2024caa}, and it behaves as a persona-level
trait, so that an off-the-shelf sceptical persona recovers most of the
reduction that targeted steering achieves \citep{chen2025persona,kelkar2026devil}.
Post-hoc remedies such as synthetic fine-tuning data reduce it
\citep{wei2023synthetic} but do not tell us what the model would conclude
if the deference were removed at inference time. We treat the persona-level
finding as licence for the intervention studied here. If sycophancy is a
stance the model adopts, then assigning a different stance is the natural
control.

\paragraph{Role assignment controls what a model says, at a cost.}
Assigning a model a character or an interest reliably moves its output
\citep{shanahan2023roleplay,li2023camel}. The literature is divided on how
much. Persona prompts explain little variance in simulated survey judgement
\citep{hu2024persona} and can degrade zero-shot reasoning
\citep{kim2025persona}, while persona-assigned models carry deep implicit
biases \citep{gupta2024bias,deshpande2023toxicity} and agreeable personas
amplify sycophancy \citep{shah2026toonice}. Multi-persona prompting within a
single model improves some tasks \citep{wang2024spp}. Closest to our design,
\citet{dietrich2026roles} assign opposed advocate roles plus a balanced seat
and a supervisor, and measure how far each advocate's output is fixed by its
role. We measure the same quantity, which we call role-lock, and find it
near total (0.97). We add the cost that the persona literature has not
priced. A role-locked seat's verdict carries almost no information about the
case, so the control and the loss of information are the same event.

\paragraph{Debate places truth at the judge, not in the debaters.}
The scalable-oversight programme puts two opposed agents in front of a judge
\citep{irving18}. Human and LLM judges become more accurate when they watch
a debate than when they consult a single advocate
\citep{michael2023debate,khan2024debating}, but \citet{kenton2024scalable}
show the gain over the judge answering alone appears only under information
asymmetry, and vanishes when the judge can see the item. Multi-agent debate
among copies of one model improves some benchmarks
\citep{du23,liang2024divergent,chan2023chateval} and fails to beat
self-consistency on others \citep{smit2024mad,wang2024rethinking}. Formal
accounts explain why. Same-model debate converges to the majority opinion
already in the room \citep{estornell2024multillm}, homogeneous agents with
uniform updates preserve expected correctness and so cannot reliably improve
it \citep{zhu2026demystifying}, and exchange is expectation-preserving in a
way that voting is not \citep{choi2025debateorvote}. Two strands use
disagreement as a signal rather than as a route to a verdict. Cross
examination between models detects factual errors \citep{cohen2023lmvslm},
and a prover--verifier game trains legibility for a checker
\citep{kirchner2024prover}. Our contribution to this lineage is to assign
opposed interests rather than opposed answers, so that neither advocate holds
the truth by construction, and then to ask two questions the lineage has
not. What do the advocates' verdicts carry, and what does their disagreement
carry. The answers are nothing and a transferable error signal.

\paragraph{Communication topology as a variable.}
Sparse debate topologies match or beat fully connected ones at lower cost
\citep{li2024sparse}, collaboration follows a scaling law over chain, tree
and graph structures \citep{qian2025macnet}, and agent systems can be
represented and optimised as graphs \citep{zhuge2024gptswarm,liu2024dylan}.
Topology governs how information and errors propagate
\citep{shen2025propagation,han2026conformity}, agents conform to a majority
\citep{zhang2024socialpsych,weng2025conformity}, and errors across models are
correlated in ways that limit what any aggregation can recover
\citep{kim2025correlated}. Heterogeneous panels outperform homogeneous ones
\citep{wang2025moa}. \citet{bertalanic2026ringelmann} define an effective
team size as a scalar order parameter for a collective and fit it across
debate, self-correction and sparse-communication conditions, and
\citet{bertalanic2026costconsensus} find isolated self-correction beats
unguided homogeneous debate. \citet{he2026minoritysentinel} ask when a
minority should overturn a majority. Across this strand the dependent
variable is task accuracy or cost. We make the structure of dissent and its
transfer to other generators the dependent variable, and we measure the
effective voter count on an embodied topology, where role assignment
supplies the mechanism that sets it to one.

\paragraph{Aggregation theory predicts the measurements.}
Condorcet's theorem requires independent voters. Correlated votes can help
or hurt majority rule \citep{ladha1992condorcet,boland1989majority}, and the
optimal rule over voters of unequal competence is a log-odds weighted
majority \citep{shapley1984optimizing,nitzan1982optimal}. With our measured
seat accuracies that rule is a dictatorship of the competent seat, which is
what we observe. \citet{golub2010naive} prove that a network converges to
the truth if and only if no node's influence dominates. Our neutral seat is
a dictator, so the collective cannot be wise, and repeated averaging can only
return the weighted prior \citep{degroot1974reaching}. Social influence
undermines crowd wisdom \citep{lorenz2011social} unless the network is
decentralised \citep{becker2017network}, and diversity is epistemically
valuable only while it is transient \citep{zollman2010epistemic}. In humans,
two heads beat one only at similar competence \citep{bahrami2010optimally},
and a few debates can outperform a large crowd \citep{navajas2018aggregated}.
The diversity-trumps-ability result \citep{hong2004groups} does not hold in
general \citep{thompson2014diversity}. Aggregating shared biases needs a
non-majoritarian rule \citep{prelec2017solution}, and judgement aggregation
faces impossibility results of its own \citep{list2002aggregating}. We read
our results as instances of these theorems rather than as surprises, and we
use the no-dominant-node condition as the design constraint any successor
topology must satisfy.

\paragraph{The argumentative theory of reasoning.}
\citet{mercier2011argumentative} and \citet{mercier2017enigma} argue that
human reasoning evolved for argument. Individuals produce biased advocacy,
and truth emerges on the evaluation side of an exchange. Groups reason well
because arguments, not confidence, move them \citep{trouche2014arguments},
while individuals are selectively lazy about their own arguments
\citep{trouche2016selective}. The theory predicts role-lock exactly. It does
not predict that advocates converge, and reviews of the evidence find the
collective half is the weaker one \citep{mercier2016predictions}. The
organisational literature reaches the same place. Assigned devil's advocacy
improves decisions less than authentic dissent \citep{nemeth2001devil},
though structured conflict beats consensus \citep{schweiger1986group,schwenk1990effects}.
Philosophically, objectivity is a property of a community's structure of
criticism rather than of any member \citep{longino1990science,longino2002fate,popper2011open}.
Our results are consistent with the theory on production and silent on
convergence, because our collective did not converge. Truth entered from
outside the graph.

\paragraph{Cooperative AI, mediation, and mechanism design.}
Cooperative AI asks how agents with divergent interests reach joint welfare
\citep{dafoe2020cooperative}. LLMs already play repeated social dilemmas,
cooperating well in prisoner's-dilemma families and poorly in coordination
games \citep{akata2025repeated}, and societies of LLM agents mostly fail to
sustain a common resource unless they communicate and reason about
long-run consequences \citep{piatti2024cooperate}. Assigned personas move
agents between cooperation and exploitation \citep{leon2026personas}.
\citet{tewolde2026coopeval} benchmark four cooperation-sustaining mechanisms
and find third-party mediation and contracting the most effective. The
formal basis is old. A mediator's correlated recommendation can reach
outcomes outside the convex hull of Nash equilibria \citep{aumann1974subjectivity}.
This is where our graph-theoretic reading meets alignment. A collective of
role-locked advocates converges to Nash, but a third node that is not locked
is exactly the mediator of correlated equilibrium and the judge of the
oversight literature. Whether such a node can pull an embodied collective
to the Pareto outcome, and whether the collective's dissent still flags the
cases where it cannot, is the question our strategic-games instrument is
built to answer.
